\begin{table}[!t]
\centering
\caption{Median performance of each DRAM assist technique.}
\label{tab:assist_summary_main}
\footnotesize
\begin{tabular}{lrrrrrrrrr}
\toprule
assist\_config & retention\_time\_s & read\_delay\_ns & write\_delay\_ns & read\_margin\_mv & total\_power\_nw & energy\_per\_access\_fj & leakage\_current\_fa & refresh\_interval\_ms & feasible\_pct \\
\midrule
baseline & 0.2865 & 5.829 & 1.224 & 39.72 & 3.040 & 28.65 & 1.509 & 301.7 & 26.33 \\
bitline\_precharge\_opt & 0.8992 & 5.706 & 1.264 & 49.58 & 3.486 & 32.54 & 1.515 & 531.8 & 25.33 \\
combined\_low\_power & 0.2812 & 5.468 & 1.359 & 67.90 & 5.784 & 52.51 & 486.4 & 151.2 & 74 \\
high\_k\_capacitor & 0.05399 & 8.256 & 1.524 & 36.31 & 4.574 & 43.31 & 19.58 & 69.27 & 17 \\
negative\_wordline & 0.1806 & 5.835 & 1.222 & 39.72 & 3.039 & 28.63 & 18.18 & 212.1 & 25.33 \\
refresh\_interval\_opt & 0.2865 & 5.829 & 1.224 & 39.72 & 3.040 & 28.65 & 1.509 & 301.7 & 26.33 \\
sense\_amp\_upsize & 0.2865 & 3.989 & 1.225 & 39.28 & 3.420 & 32.66 & 1.508 & 301.7 & 28.00 \\
wordline\_boost & 1.280 & 5.266 & 0.8842 & 61.93 & 3.233 & 29.26 & 2.674 & 675.5 & 76.33 \\
\bottomrule
\end{tabular}
\\[2pt] \footnotesize Paired study: every technique is applied to the identical set of 300 base designs and simulated in NGSpice.
\end{table}
